\documentclass[13pt,a4paper]{article}

% Paquetes para idioma, codificación y márgenes
\usepackage[utf8]{inputenc}
\usepackage[spanish]{babel}
\usepackage[margin=2.5cm]{geometry}
\usepackage{enumitem} % Para mejorar el espaciado de las listas

% Ajustes de párrafo
\setlength{\parindent}{0pt}
\setlength{\parskip}{1em}

\begin{document}

\begin{flushright}
\textbf{Fecha:} \today
\end{flushright}

\vspace{0.5cm}

\textbf{Asunto:} Metodología de trabajo y recursos tecnológicos para Dibujo Técnico Asistido

\vspace{0.5cm}

Estimadas familias:

Les doy la bienvenida al ciclo lectivo. Me comunico mediante la presente para informarles cómo trabajaremos este año en el espacio de \textbf{Dibujo Técnico Asistido}.

En esta etapa, los alumnos darán un paso fundamental en su formación técnica: pasarán del uso del tablero de dibujo tradicional al diseño digital, utilizando herramientas de software de nivel profesional como \textbf{FreeCAD} (diseño 3D) y \textbf{KiCad} (diseño electrónico).

Para realizar estas prácticas, la herramienta de estudio indispensable es una computadora portátil (notebook o netbook) con sistema operativo Windows o Linux y al menos 2\,GB de memoria RAM. \textbf{Es importante aclarar que los teléfonos celulares y las tablets, sin importar su gama o tecnología, no son aptos ni compatibles para instalar o manejar estos programas de diseño técnico.}

Como es de conocimiento general, al no estar vigentes los programas estatales de provisión de equipos, la escuela cuenta actualmente con una cantidad muy limitada de netbooks de antigua data (con 1\,GB de RAM). Si bien estas máquinas logran ejecutar los programas, lo hacen con una marcada lentitud frente a las exigencias del software actual.

Por este motivo:
\begin{itemize}[itemsep=2pt, parsep=0pt]
    \item \textbf{Recomendamos fuertemente} que aquellos alumnos que dispongan de una notebook o netbook en casa la traigan a las clases para usarla como su estación de trabajo personal.
    \item Los alumnos que no cuenten con equipo propio \textbf{trabajarán de forma colaborativa} compartiendo las computadoras disponibles en la escuela.
\end{itemize}

Desde mi rol, haré todo lo posible para que todos aprendan y completen sus actividades. Sin embargo, es mi deber informar que el trabajo compartido y las limitaciones técnicas de los equipos escolares pueden hacer que los tiempos de práctica individual y el avance general en los trabajos prácticos sean más lentos de lo ideal.

Agradezco de antemano su comprensión y el apoyo desde casa para con la educación técnica de los alumnos. Quedo a su entera disposición por cualquier duda o consulta.

\vspace{2cm}

\begin{flushleft}
Dibujo Técnico Asistido
\end{flushleft}

\end{document}
